\documentclass[12pt,addpoints]{exam}

% ---------- Packages ----------
\usepackage{amsmath,amssymb}
\usepackage{svg}
\usepackage{geometry}
\usepackage[dvipsnames]{xcolor}

% Enable solution printing
\printanswers

\geometry{margin=1in}

\definecolor{light1}{HTML}{f1f2f6}
\definecolor{blue1}{HTML}{30336b}

% Enable highlight-boxes
\usepackage[most]{tcolorbox}

%textmarker style from colorbox doc
\tcbset{textmarker/.style={%
        enhanced,
        parbox=false,boxrule=0mm,boxsep=0mm,arc=0mm,
        outer arc=0mm,left=6mm,right=3mm,top=7pt,bottom=7pt,
        toptitle=1mm,bottomtitle=1mm,oversize}}

\newtcolorbox{lightBox}{textmarker,
    borderline west={6pt}{0pt}{light1},
    colback=light1!10!white}

\newtcolorbox{blueBox}{textmarker,
    borderline west={6pt}{0pt}{blue1},
    colback=blue1!10!white}

\newcommand{\bluebox}[1]{\begin{blueBox} #1 \end{blueBox}}

% ---------- Solution Box Styling ----------
\newtcolorbox{solutionstyle}{
    colback=blue!5!white,
    colframe=blue!75!black,
    title=\textbf{Solution},
    arc=2mm,
    boxrule=0.5pt
}

% Override the default exam solution environment
\renewenvironment{solution}{\begin{solutionstyle}}{\end{solutionstyle}}

% Placeholder boxes (unused when \printanswers is on, but kept for reference)
\newcommand{\answerboxS}{\begin{lightBox} \vspace{3cm} \end{lightBox}}
\newcommand{\answerboxM}{\begin{lightBox} \vspace{5cm} \end{lightBox}}
\newcommand{\answerboxL}{\begin{lightBox} \vspace{7cm} \end{lightBox}}

% ---------- Header / Footer ----------
\pagestyle{headandfoot}
\runningheadrule
\runningheader{}{Numerical Programming I - Solution Key}{Page \thepage\ of \numpages}
\firstpagefooter{}{}{}
\runningfooter{}{}{}

% ---------- Point Boxes ----------

\boxedpoints
\pointsinmargin

\begin{document}

{\noindent\Huge\bf  \\[0.5\baselineskip] {\fontfamily{cmr}\selectfont  Exam -- Reference Solution}}\\[2\baselineskip]
{ {\bf \fontfamily{cmr}\selectfont Numerical Programming I}\\
{\textit{\fontfamily{cmr}\selectfont     February 6, 2026}}}
\\[1.4\baselineskip] 

% ---------- Student Info ----------
\makebox[0.75\textwidth]{Name:\enspace\hrulefill}

\vspace{2cm}

% ---------- Grading Table ----------
\begin{center}
\gradetable[h][questions]
\end{center}

\pagebreak

% =========================================================
%                    EXAM BEGINS
% =========================================================

\begin{questions}

% =========================================================
% MULTIPLE CHOICE SECTION
% =========================================================
\section*{Mixed Short Questions}

\question
\begin{parts}
\part[1] What does $(1.0 + 10^{-17}) - 1.0$ evaluate to in double-precision floating point arithmetic (with machine precision $\epsilon_\text{mach} \sim 10^{-16}$)?
\begin{checkboxes}
  \choice $10^{-17}$ is correctly calculated.
  \choice $10^{-17}$ cannot be represented in double-precision, the expression results in undefined behavior.
  \CorrectChoice The expression evaluates to zero, the mantissa of $1.0$ does not have enough digits to capture $10^{-17}$.
\end{checkboxes}
\begin{solution}
    Since $10^{-17} < \epsilon_{\text{mach}}/2$ relative to $1.0$, the addition $1.0 + 10^{-17}$ results in exactly $1.0$ due to rounding (machine epsilon is the smallest number such that $1+\epsilon \neq 1$). Subtracting $1.0$ yields $0.0$.
\end{solution}

\part[2] Consider a differentiable function $f: \mathbb{R} \to \mathbb{R}$. We have written a program
which approximates $f'(x)$ using the forward difference $\frac{f(x+\Delta x) - f(x)}{\Delta x}$. 
Below, the absolute error $\left| f'(x) - \frac{f(x+\Delta x) - f(x)}{\Delta x}\right|$ over $\Delta x$ is plotted. Mark the
region where the floating point error is dominant and the region where the truncation error is dominant.

\includesvg[width=0.7\textwidth]{figures/forward_error.svg}

\begin{solution}
    \textbf{Left Region (jagged/noisy lines):} Floating point cancellation error is dominant (as $\Delta x \to 0$).\\
    \textbf{Right Region (straight slope 1 line):} Truncation error is dominant (mathematical approximation error).
\end{solution}

\part[2] Does gradient descent always converge to the \textbf{global} minimum? Answer with yes or no and give a short justification.
\begin{solution}
    \textbf{No.} Gradient descent follows the steepest descent direction locally. It will converge to a \textbf{local} minimum depending on the initialization. It only guarantees global convergence for strictly convex functions.
\end{solution}
\end{parts}

\pagebreak

% =========================================================
% LONG FORM QUESTIONS
% =========================================================
\section*{Second order derivative approximation}

\question
For a twice differentiable function $f: \mathbb{R} \to \mathbb{R}$
we would like to find a finite difference approximation for the
second derivative $f''(x)$.

\begin{parts}

\part[2] Approximate $f(x+\Delta x)$ and $f(x - \Delta x)$ by Taylor
expansion up to (including) third order around $x$.
\begin{solution}
    The Taylor expansions are:
    \begin{align*}
        f(x+\Delta x) &= f(x) + f'(x)\Delta x + \frac{1}{2}f''(x)\Delta x^2 + \frac{1}{6}f'''(x)\Delta x^3 + \mathcal{O}(\Delta x^4) \\
        f(x-\Delta x) &= f(x) - f'(x)\Delta x + \frac{1}{2}f''(x)\Delta x^2 - \frac{1}{6}f'''(x)\Delta x^3 + \mathcal{O}(\Delta x^4)
    \end{align*}
\end{solution}

\part[2] Using these approximations find an approximation for $f''(x)$
based on $f(x)$, $f(x+\Delta x)$ and $f(x-\Delta x)$. What is the order of the error term?
\textit{Hint:} You need to combine both
Taylor expansions such that the first and third order derivatives cancel.
\begin{solution}
    Summing the two equations from part (a):
    \begin{align*}
        f(x+\Delta x) + f(x-\Delta x) &= 2f(x) + f''(x)\Delta x^2 + \mathcal{O}(\Delta x^4)
    \end{align*}
    (Note that the odd terms $f'$ and $f'''$ cancel out). \\
    Solving for $f''(x)$:
    \begin{align*}
        f''(x)\Delta x^2 &= f(x+\Delta x) - 2f(x) + f(x-\Delta x) + \mathcal{O}(\Delta x^4) \\
        f''(x) &= \frac{f(x+\Delta x) - 2f(x) + f(x-\Delta x)}{\Delta x^2} + \mathcal{O}(\Delta x^2)
    \end{align*}
    The order of the error term is $\mathcal{O}(\Delta x^2)$ (second order).
\end{solution}

\part[1] For $f(x) = x^2$ and $\Delta x = 0.1$ approximate $f''(0)$ and compare
your approximation to the analytical value.
\begin{solution}
    \textbf{Analytical:} $f(x)=x^2 \implies f'(x)=2x \implies f''(x)=2$. Thus $f''(0)=2$. \\
    \textbf{Approximation:} Using the formula derived above at $x=0$:
    \begin{align*}
        f''(0) &\approx \frac{f(0.1) - 2f(0) + f(-0.1)}{(0.1)^2} \\
               &= \frac{0.01 - 0 + 0.01}{0.01} = \frac{0.02}{0.01} = 2
    \end{align*}
    The approximation matches the analytical value exactly (because the error term depends on $f^{(4)}$, which is 0 for a quadratic).
\end{solution}

\end{parts}

\pagebreak

\section*{Convergence of Newton’s root finding method}

\question

In the lecture we have stated:

\bluebox{If $g$ is continuously differentiable in a neighborhood
around a root $x^\star$ with nonzero derivative at $x^\star$, 
then there exists a neighborhood around $x^\star$ where Newton's 
iteration converges to $x^\star$ for all starting values.}

Now consider the function

\begin{equation}
    g(x) = x(1+x^2)^{-\frac{1}{2}}
\end{equation}

\begin{parts}
\part[1] What is the root $x^\star$ of $g$ where $g(x^\star) = 0$?

\begin{solution}
    $g(x) = \frac{x}{\sqrt{1+x^2}}$. Since the denominator $\sqrt{1+x^2} \geq 1$ never vanishes, the only solution is the numerator being zero.
    $$ x^\star = 0 $$
\end{solution}

\part[2] Show that the derivative of $g$ is $g'(x) = (1+x^2)^{-\frac{3}{2}}$. \textit{Hint:} Calculate the
derivative by product rule and factor out $(1+x^2)^{-\frac{3}{2}}$.

\begin{solution}
    Using the product rule on $x \cdot (1+x^2)^{-1/2}$:
    \begin{align*}
        g'(x) &= 1 \cdot (1+x^2)^{-\frac{1}{2}} + x \cdot \left(-\frac{1}{2}\right)(1+x^2)^{-\frac{3}{2}} \cdot (2x) \\
              &= (1+x^2)^{-\frac{1}{2}} - x^2(1+x^2)^{-\frac{3}{2}} \\
        \text{Factor out } (1+x^2)^{-\frac{3}{2}}: \\
              &= (1+x^2)^{-\frac{3}{2}} \left[ (1+x^2)^1 - x^2 \right] \\
              &= (1+x^2)^{-\frac{3}{2}} \cdot 1 = (1+x^2)^{-\frac{3}{2}}
    \end{align*}
\end{solution}

\part[1] Calculate $g'(x^\star)$. Is there a neighborhood around $x^\star$ where Newton's 
iteration converges to $x^\star$ for all starting values inside that neighborhood?

\begin{solution}
    $$ g'(0) = (1+0^2)^{-\frac{3}{2}} = 1 $$
    Since $g'(x^\star) = 1 \neq 0$, the condition of the theorem holds. \textbf{Yes}, there exists such a neighborhood.
\end{solution}

\part[2] Show that the Newton step
\begin{equation}
    x_{k+1} = x_{k} - \frac{g(x_k)}{g'(x_k)}
\end{equation}
is given by $x_{k+1} = - x_k^3$.

\begin{solution}
    \begin{align*}
        \frac{g(x)}{g'(x)} &= \frac{x(1+x^2)^{-\frac{1}{2}}}{(1+x^2)^{-\frac{3}{2}}} \\
                           &= x (1+x^2)^{-\frac{1}{2}} (1+x^2)^{\frac{3}{2}} \\
                           &= x (1+x^2)^{1} = x + x^3
    \end{align*}
    Therefore:
    $$ x_{k+1} = x_k - (x_k + x_k^3) = -x_k^3 $$
\end{solution}

\part[3] Consider the following cases for the initialization of 
the Newton iteration
\begin{itemize}
    \item $x_0 = 0.5$
    \item $x_0 = -1.0$
    \item $x_0 = 2.0$
\end{itemize}
and analyze if the iteration converges, diverges or oscillates.

\begin{solution}
    Using $x_{k+1} = -x_k^3$:
    \begin{itemize}
        \item $x_0 = 0.5$: $x_1 = -0.125$, $x_2 = -(-0.125)^3 \approx 0.002$. Values get smaller rapidly. \textbf{Converges} to 0.
        \item $x_0 = -1.0$: $x_1 = -(-1)^3 = 1.0$, $x_2 = -(1)^3 = -1.0$. \textbf{Oscillates} between -1 and 1.
        \item $x_0 = 2.0$: $x_1 = -(2)^3 = -8$, $x_2 = -(-8)^3 = 512$. Values grow rapidly. \textbf{Diverges}.
    \end{itemize}
\end{solution}

\end{parts}

\pagebreak

\section*{\textit{(Bonus)} Numerical integration by hand}

\bonusquestion

Consider the function $f: \mathbb{R} \to \mathbb{R}$ plotted below.

\includesvg[width=0.7\textwidth]{figures/integration_by_hand.svg}

We want to approximate the integral

\begin{equation}
    \int_{0}^{5} f(x) \, dx
\end{equation}

Here, assume the subdivision

\begin{equation}
    \int_{0}^{5} f(x) \, dx = \sum_{i=0}^{N - 1} \int_{x_i}^{x_{i+1}} f(x) \, dx
\end{equation}

with $N = 5$ and $x_0 = 0.0, x_1 = 1.0, x_2 = 2.0, x_3 = 3.0, x_4 = 4.0, x_5 = 5.0$.

And consider the following \textit{quadrature rules}:

\begin{itemize}
    \item midpoint rule:
    \begin{equation}
        \int_{x_i}^{x_{i+1}} f(x) \, dx = \underbrace{(x_{i+1} - x_i) \, f\left( \frac{x_{i} + x_{i+1}}{2} \right)}_{=:M_i} + \mathcal{O}((x_{i+1} - x_i)^3)
    \end{equation}
    \item trapezoidal rule:
    \begin{equation}
        \int_{x_i}^{x_{i+1}} f(x) \, dx = \underbrace{(x_{i+1} - x_i) \, \left(\frac{f(x_i) + f(x_{i+1})}{2}\right)}_{=:T_i} + \mathcal{O}((x_{i+1} - x_i)^3)
    \end{equation}
\end{itemize}

\begin{parts}
    \bonuspart[2] Approximate $\int_{0}^{5} f(x) \, dx$ using the midpoint rule as given above.
    \begin{solution}
        For the Midpoint rule with $N=5$, we need values at $x = 0.5, 1.5, 2.5, 3.5, 4.5$. Reading the blue dots on the graph:
        \begin{itemize}
            \item $f(0.5) = 3.0$
            \item $f(1.5) = 2.5$
            \item $f(2.5) = 1.5$
            \item $f(3.5) = 2.0$
            \item $f(4.5) = 1.5$
        \end{itemize}
        The midpoint approximation yields
        \begin{equation}
            M = \Delta x \sum f(midpoints) = 1.0 \cdot (3.0 + 2.5 + 1.5 + 2.0 + 1.5) = \mathbf{10.5}
        \end{equation}
    \end{solution}

    \bonuspart[2] Approximate $\int_{0}^{5} f(x) \, dx$ using the trapezoidal rule as given above.
    \begin{solution}
        For the Trapezoidal rule, we need values at integers $x=0, 1, 2, 3, 4, 5$. Reading the graph:
        \begin{itemize}
            \item $f(0)=2.0$
            \item $f(1)=3.0$
            \item $f(2)=2.0$
            \item $f(3)=1.5$
            \item $f(4)=2.0$
            \item $f(5)=1.0$
        \end{itemize}
        The trapezoidal approximation yields
        \begin{equation}
            \begin{aligned}
                T &= \Delta x \sum_{i=0}^{4} \frac{f(x_i) + f(x_{i+1})}{2} \\
                &= \Delta x \left[ \frac{f(0)+f(5)}{2} + \sum_{i=1}^{4} f(x_i) \right] \\
                &= 1.0 \cdot \left[ \frac{2.0 + 1.0}{2} + (3.0 + 2.0 + 1.5 + 2.0) \right] \\
                &= \mathbf{10.0}
            \end{aligned}
        \end{equation}
    \end{solution}
\end{parts}


\end{questions}

\end{document}
